\documentclass[twoside,11pt]{article}

\usepackage{jmlr2e}
\usepackage{amsmath}
\usepackage{microtype}
\usepackage{booktabs}
\usepackage{tabularx}
\usepackage{array}
\usepackage{enumitem}
\usepackage{lastpage}

\jmlrheading{0}{2026}{1-\pageref{LastPage}}{Submitted 4/2026}{Preprint}{PAPER-C-0001}{Jonathan Herrera-Vasquez}
\ShortHeadings{From Fidelity to Semantics}{Herrera-Vasquez}
\firstpageno{1}

\begin{document}

\title{From Fidelity to Semantics:\\
A Taxonomy and Survey of Evaluation Metrics for Model-Agnostic Explainability}

\author{
\name Jonathan Herrera-Vasquez \email jonnabio@gmail.com \\
\addr UNADE
}

\editor{}

\maketitle

\begin{abstract}
Evaluation remains the least consolidated part of explainable artificial
intelligence (XAI). While model-agnostic post-hoc methods such as LIME,
SHAP, Anchors, and DiCE are widely used, the criteria used to judge them
vary sharply across papers, toolkits, and application domains. Some
studies emphasize faithfulness proxies, others prioritize stability,
sparsity, or human plausibility, and newer work introduces large language
models (LLMs) as semantic evaluators without a unified framework for
locating those judgments within the broader evaluation landscape. This
paper presents a structured survey and taxonomy of evaluation metrics for
model-agnostic explainability, focusing on how technical proxy metrics,
benchmark-grounded metrics, human evaluation, and LLM-based semantic
evaluation relate to one another. Using a coded literature matrix derived
from the thesis review corpus, we organize XAI evaluation along four
axes: evaluation target, evidence source, quality property, and task
context. The resulting taxonomy clarifies where current metrics overlap,
where they leave construct gaps, and why semantic evaluation should be
treated as a complementary layer rather than a replacement for
robustness- and fidelity-oriented measurement. The paper closes by
mapping the taxonomy back to the thesis framework, motivating the
benchmark design of Paper B and the evaluator-validation agenda of Paper
A.
\end{abstract}

\begin{keywords}
explainable ai, evaluation metrics, survey, taxonomy, semantic
evaluation, llm-as-a-judge
\end{keywords}

\section{Introduction}

Progress in explainable artificial intelligence depends not only on
generating explanations, but on evaluating whether those explanations
are useful, faithful, stable, and meaningful. Yet the literature does
not converge on a single definition of explanation quality. Surveys
repeatedly note that XAI evaluation is fragmented across proxy metrics,
qualitative user studies, and domain-specific desiderata
\citep{adadi2018xai,arrieta2020xai,kadir2023metrics,canha2025functionally}.
As a result, comparisons between explainers often depend as much on the
chosen metric family as on the methods being compared.

This fragmentation is especially visible in model-agnostic post-hoc
explainability. Faithfulness-style metrics remain common because they
are relatively easy to compute, but they do not fully capture whether an
explanation is understandable, semantically aligned with domain
knowledge, or decision-useful in practice
\citep{ali2023trustworthy,longo2024manifesto}. Multi-metric work argues
that stability, sparsity, and computational cost are also part of
explanation quality, and recent critical examinations show that
different metrics can reveal complementary rather than redundant
behavior \citep{pawlicki2024multiple,zheng2025ffidelity}.

At the same time, a new line of work uses LLMs as evaluators of
explanation quality or human-centered criteria. This is promising
because LLMs can score natural-language justifications, compare
competing explanations, and operationalize semantic rubrics at scale.
However, research on LLM-as-a-judge also warns that proxy judges can
drift from human judgments, inherit style biases, and compress distinct
quality dimensions into superficially coherent scores
\citep{zheng2023judging,gu2024llmjudge,wu2025userperceptions}. For XAI,
this raises a conceptual problem before it raises an empirical one:
where should semantic evaluation sit relative to fidelity, robustness,
and other established metric families?

This paper addresses that question through a structured survey and
taxonomy for model-agnostic XAI evaluation. Rather than proposing a new
explainer or a new benchmark, the paper organizes the metric space
needed to interpret the rest of the thesis. In the current dissertation
plan, Paper B provides the quantitative benchmark layer and Paper A
validates LLM-based semantic evaluation against human judgments. Paper C
supplies the conceptual bridge between them.

Our contributions are:
\begin{enumerate}[leftmargin=*]
\item a taxonomy of XAI evaluation metrics structured by evaluation
target, evidence source, quality property, and task context;
\item a comparative synthesis of technical proxy metrics,
benchmark-grounded metrics, human evaluation, and LLM-based semantic
evaluation;
\item a gap analysis showing where current metrics fail to capture
explanation meaning or user-facing adequacy; and
\item a thesis-facing integration layer that motivates why Paper B needs
a multi-metric benchmark and why Paper A needs explicit evaluator
validation.
\end{enumerate}

Scope is intentionally limited to model-agnostic post-hoc
explainability, with emphasis on tabular settings but attention to
transfer issues across image, text, time series, and graph contexts.
The paper is a structured narrative survey with explicit coding
dimensions, not a finalized systematic review or meta-analysis.

\section{Background and Definitions}

To keep the taxonomy precise, several terms that are often conflated in
XAI writing need to be separated.

First, intrinsic explainability and post-hoc explainability should not
be evaluated as if they pose the same problem. Intrinsic models are
designed to be interpretable by construction, whereas post-hoc methods
attempt to explain black-box behavior after model training
\citep{rudin2019stop,retzlaff2024posthoc}. The present survey focuses on
the latter because Papers A and B evaluate post-hoc model-agnostic
explainers.

Second, explanation quality is not identical to model quality. A model
may be accurate while producing unstable or semantically weak
explanations; conversely, a model can support concise explanations
without necessarily being more accurate. Evaluation therefore needs to
specify whether the object under study is the explanation artifact, the
explainer method, the predictive model, or the user-facing decision
process.

Third, local and global explanations require different expectations.
Local attributions or counterfactuals are often judged through
faithfulness, sparsity, or actionability, while global summaries may be
judged by coverage, consistency, or conceptual coherence. Many reported
metric disagreements arise because these explanatory levels are mixed
without explicit task framing.

Finally, semantic evaluation should be separated from purely structural
proxy metrics. A metric such as perturbation faithfulness evaluates
whether feature rankings track model sensitivity under a chosen
intervention scheme. A semantic judgment evaluates whether an
explanation ``makes sense'' under domain, causal, or linguistic
criteria. Both matter, but they are not interchangeable.

\section{Review Method and Coding Protocol}

This prototype uses a structured narrative review protocol anchored in
the repository's literature matrix,
\path{docs/reports/literature_review_methodology_matrix.md}. The current
coded set contains 25 papers spanning metric taxonomies, benchmark
toolkits, faithfulness theory, modality-specific benchmarks, and
LLM-as-a-judge research. The paper set was originally assembled to
support thesis methodology design rather than to satisfy a pre-registered
systematic review workflow, so the present manuscript treats it as a
transparent review corpus rather than as an exhaustive bibliometric
sample.

The review objective is to answer four organizing questions:
\begin{enumerate}[leftmargin=*]
\item What kinds of objects are current XAI metrics actually evaluating?
\item What evidence sources do they rely on?
\item Which quality properties do they operationalize?
\item Where do semantic or human-centered judgments become necessary?
\end{enumerate}

Each source in the coded corpus is assigned to one or more thematic
clusters:
\begin{itemize}[leftmargin=*]
\item metric taxonomies and surveys;
\item benchmark and toolkit papers;
\item faithfulness and robustness studies;
\item modality-specific or domain-specific evaluations; and
\item LLM-as-a-judge and proxy-human-evaluation papers.
\end{itemize}

For this prototype, inclusion is defined pragmatically: a source must
contribute either a formal evaluation construct, a benchmark design
principle, a warning about metric misuse, or an evaluation paradigm
directly relevant to post-hoc model-agnostic explanations. Exclusion
applies to papers that focus solely on generating explanations without
materially informing evaluation design.

\begin{table}[t]
\centering
\caption{Coding dimensions used to build the Paper C taxonomy.}
\begin{tabularx}{\linewidth}{l X}
\toprule
\textbf{Dimension} & \textbf{Operational question} \\
\midrule
Evaluation target & What is being judged: explanation artifact, explainer, model behavior, or user/task outcome? \\
Evidence source & What kind of evidence supports the judgment: perturbation proxy, benchmark ground truth, human rating, user task, or LLM judge? \\
Quality property & Which property is being claimed: fidelity, stability, sparsity, plausibility, recourse quality, semantic alignment, and related constructs? \\
Task context & In what modality or domain is the metric assumed to be meaningful? \\
\bottomrule
\end{tabularx}
\label{tab:coding}
\end{table}

This protocol is intentionally lighter than PRISMA-style systematic
review practice. The goal of the prototype is conceptual organization
and thesis integration, not final exhaustive coverage claims.

\section{Taxonomy of XAI Evaluation Metrics}

The central claim of this paper is that XAI evaluation can be clarified
by crossing four axes: evaluation target, evidence source, quality
property, and task context. Metrics that appear to answer the same
question often do not, because they attach to different objects, use
different evidence, or imply different validity conditions.

\subsection{Taxonomy by Evaluation Target}

Many metric disagreements start with a hidden mismatch about what is
being evaluated.
\begin{itemize}[leftmargin=*]
\item \textbf{Explanation artifact metrics} evaluate the produced
explanation itself, such as attribution sparsity, local fidelity, or
semantic clarity.
\item \textbf{Explainer method metrics} evaluate the behavior of the
explainer across cases, such as stability, reproducibility, or
computational cost.
\item \textbf{Model-behavior metrics} evaluate whether explanations
track model behavior under perturbations or alternative decision
regions.
\item \textbf{User/task outcome metrics} evaluate whether explanations
improve decision performance, trust calibration, or actionability.
\end{itemize}

This distinction matters because one metric family cannot automatically
stand in for another. A sparse explanation can still be semantically
misleading, and a faithful attribution can still be unusable in a
decision-support workflow.

\subsection{Taxonomy by Evidence Source}

The second axis concerns the evidence used to justify the evaluation:
\begin{itemize}[leftmargin=*]
\item \textbf{Proxy or perturbation metrics} infer quality from
numerical sensitivity behavior.
\item \textbf{Benchmark-grounded metrics} compare explanations to
synthetic or domain-defined ground truth.
\item \textbf{Human expert judgments} evaluate plausibility,
correctness, or domain alignment.
\item \textbf{End-user studies} assess comprehension, trust
calibration, decision quality, or workload.
\item \textbf{LLM judges} provide scalable semantic scoring using
rubric-guided language-model evaluation.
\end{itemize}

These evidence sources are not equally strong for every construct.
Proxy metrics are scalable and reproducible but may under-capture
meaning. Human judgments are richer but expensive and variable. LLM
judges sit in between: they are scalable like proxies but semantically
expressive like human reviews, which is precisely why their validity
must be established rather than assumed
\citep{gu2024llmjudge,wu2025userperceptions,zhou2025medthink}.

\subsection{Taxonomy by Quality Property}

The literature repeatedly returns to a familiar set of quality
properties, but not always with consistent definitions.
\begin{itemize}[leftmargin=*]
\item \textbf{Fidelity / faithfulness:} whether the explanation tracks
how the model actually behaves under feature interventions or masking
\citep{ribeiro2016why,lundberg2017unified,zheng2025ffidelity}.
\item \textbf{Stability / robustness:} whether similar inputs yield
similar explanations under perturbation or reruns
\citep{pawlicki2024multiple,canha2025functionally}.
\item \textbf{Sparsity / parsimony:} whether the explanation is concise
enough to be interpretable without unnecessary feature load.
\item \textbf{Plausibility:} whether the explanation aligns with domain
expectations or expert reasoning.
\item \textbf{Counterfactual or recourse quality:} whether
explanation-linked interventions are feasible, minimal, diverse, or
actionable \citep{mothilal2020dice}.
\item \textbf{Semantic alignment:} whether an explanation captures
meaningful causal or conceptual structure, even when pure proxy metrics
are agnostic to that structure.
\end{itemize}

The key observation is that these properties are only partly
commensurable. A method can score strongly on fidelity while remaining
weak on sparsity or semantic adequacy. Multi-metric reporting is
therefore not a luxury; it is often the minimum needed to make XAI
claims interpretable.

\subsection{Taxonomy by Task Context}

Metric validity also depends on context. Tabular benchmarks often make
strong assumptions about feature independence, perturbation realism, and
local neighborhoods. These assumptions can fail in image, text, graph,
and time-series settings, where perturbations may create
out-of-distribution examples or semantically invalid transformations
\citep{agarwal2023gnn_eval,proszewska2025bxaic}. This does not make
cross-modal benchmarks impossible, but it does mean that metric
transport cannot be assumed.

\begin{table}[t]
\centering
\caption{Core metric families in model-agnostic XAI evaluation.}
\footnotesize
\begin{tabularx}{\linewidth}{l X X}
\toprule
\textbf{Metric family} & \textbf{What it captures well} & \textbf{What it leaves unresolved} \\
\midrule
Fidelity / faithfulness & Whether explanations track model behavior under interventions or masking. & Whether the explanation is meaningful, plausible, or useful outside the proxy setup. \\
Stability / robustness & Whether explanations remain consistent across perturbations or repeated runs. & Whether stable explanations are actually correct or semantically adequate. \\
Sparsity / parsimony & Whether explanations are concise and low in feature load. & Whether concision hides necessary nuance or context. \\
Benchmark-grounded correctness & Whether explanations recover synthetic or domain-defined ground truth. & Whether the available ground truth is realistic, complete, or transferable. \\
Human expert review & Whether explanations align with domain expectations and expert reasoning. & Whether those judgments are scalable, standardized, and reproducible. \\
User-centered evaluation & Whether explanations improve understanding, decisions, or trust calibration in practice. & Whether findings generalize across tasks, populations, and study designs. \\
LLM-based semantic evaluation & Whether rubric-based semantic judgments can be produced quickly and at scale. & Whether those judgments agree with humans and avoid proxy-judge bias. \\
\bottomrule
\end{tabularx}
\label{tab:taxonomy}
\end{table}

Table~\ref{tab:taxonomy} illustrates the practical consequence of the
taxonomy: each metric family gains strength by focusing on one kind of
evidence and one kind of target, but each also leaves predictable blind
spots. No single family can therefore exhaust explanation quality on its
own.

\section{Comparative Synthesis and Open Gaps}

The taxonomy exposes three recurring gaps in current XAI evaluation
practice.

\subsection{Gap 1: Proxy Metrics Are Necessary but Not Sufficient}

Faithfulness, stability, and sparsity remain essential because they
create reproducible, quantitative anchors for benchmark comparison.
Papers B and related benchmark work depend on these anchors for
statistical comparison across explainers
\citep{agarwal2022openxai,hedstrom2023quantus,canha2025functionally}.
However, these metrics operate on structural or behavioral proxies. They
do not directly answer whether an explanation reflects domain logic,
causal relevance, or user-meaningful interpretation.

\subsection{Gap 2: Human-Grounded Constructs Are Underspecified}

The literature frequently invokes interpretability, clarity,
plausibility, and usefulness, but these are often
under-operationalized. When human studies exist, their tasks, rubrics,
and participant populations vary so widely that direct comparison is
difficult \citep{ali2023trustworthy,longo2024manifesto}. This weakens
cumulative progress because the same label may refer to different
constructs across papers.

\subsection{Gap 3: Semantic Evaluation Is Rising Faster than Its Validation Base}

LLM-based judging offers a practical way to score explanation quality at
scale, especially when explanations or rubrics are textual. But the
broader LLM-as-a-judge literature shows that proxy judges can display
position bias, verbosity bias, or agreement illusions
\citep{zheng2023judging,gu2024llmjudge}. In XAI, that means semantic
evaluator outputs should be treated as calibrated proxies rather than as
self-validating replacements for expert review. This gap directly
motivates Paper A.

\begin{figure}[t]
\centering
\fbox{\begin{minipage}{0.94\linewidth}
\centering
\textbf{Evaluation continuum for post-hoc XAI}\\[4pt]
Proxy metrics $\rightarrow$ benchmark-grounded checks $\rightarrow$ expert judgment $\rightarrow$ user-centered evaluation\\
\vspace{4pt}
\footnotesize LLM-based semantic evaluation occupies an intermediate role:
more expressive than proxy metrics, more scalable than human review, but
not self-validating.
\end{minipage}}
\caption{Conceptual placement of semantic evaluation within the broader XAI evaluation landscape.}
\label{fig:continuum}
\end{figure}

Figure~\ref{fig:continuum} summarizes the evaluation continuum emphasized
by this survey. The synthesis therefore argues for layered evaluation.
Proxy metrics should anchor reproducibility and comparative benchmarking.
Human or semantic layers should be added when the claim concerns
meaning, plausibility, or causal adequacy. The mistake is not using one
family or the other; it is allowing any single family to monopolize the
definition of explanation quality.

\section{Implications for the Thesis Framework}

The value of Paper C is not only descriptive. It directly organizes the
thesis architecture.

First, the taxonomy justifies the multi-metric benchmark design of Paper
B. If explanation quality spans fidelity, stability, sparsity, and
computational burden, then a benchmark based on one endpoint would
understate meaningful trade-offs. Paper B therefore occupies the
proxy-and-robustness layer of the taxonomy.

Second, the taxonomy explains why Paper A is necessary. Once the thesis
claims to measure causal alignment, counterfactual sensitivity, or
semantic adequacy through LLM evaluators, those evaluators become
measurement instruments rather than convenience tools. Their agreement
with human experts must therefore be validated. Paper A occupies the
semantic-evaluator validation layer of the taxonomy.

Third, the survey suggests a thesis-wide minimal evaluation stack:
\begin{enumerate}[leftmargin=*]
\item technical proxy metrics for fidelity, stability, sparsity, and
cost;
\item benchmark reporting with explicit scope and artifact
traceability;
\item semantic evaluation for constructs not captured by numerical
proxies; and
\item human validation whenever semantic claims become confirmatory
rather than exploratory.
\end{enumerate}

This layered view also clarifies claim discipline. Benchmark evidence
can support comparative technical claims. Semantic evaluator outputs can
support exploratory interpretation until calibrated. Human agreement
evidence can then convert those semantic constructs into stronger
confirmatory claims.

\section{Limitations}

This prototype has several boundaries.
\begin{itemize}[leftmargin=*]
\item It is a structured narrative survey rather than a finalized
systematic review. The coded corpus is explicit and traceable, but it is
not yet a complete database-driven search with deduplication and flow
reporting.
\item The paper emphasizes model-agnostic post-hoc XAI because that is
the thesis focus. Intrinsic interpretability, representation-level
interpretability, and some modality-specific traditions are only covered
insofar as they affect metric transfer or scope boundaries.
\item The paper argues for semantic evaluation as a necessary layer but
does not itself validate LLM judges. That validation belongs to Paper A.
\item The taxonomy organizes construct space, not effect sizes. It
clarifies what should be measured and why, but it does not estimate a
pooled numerical ranking across studies.
\end{itemize}

\section{Conclusion}

XAI evaluation is not one problem but a family of related measurement
problems. The literature has developed useful metric families for
fidelity, robustness, sparsity, and benchmark reproducibility, yet
these families do not fully capture semantic adequacy or human-centered
meaning. By organizing model-agnostic post-hoc evaluation along the axes
of evaluation target, evidence source, quality property, and task
context, this paper clarifies why metric disagreement is common and why
layered evaluation is necessary.

The main practical implication is simple: benchmark metrics and semantic
evaluation should be treated as complementary rather than competing
forms of evidence. Within the thesis, this means Paper B should provide
the quantitative comparative baseline, while Paper A should validate the
semantic evaluator layer that benchmark metrics alone cannot justify.
Paper C provides the conceptual map that makes those two efforts
cohere.

\acks{This work received no external funding. The study was self-funded
by the author, and the author reports no competing interests. This draft
was prepared from repository materials available in April 2026.}

\clearpage
\phantomsection
\pdfbookmark[1]{References}{references}
\begin{thebibliography}{99}

\bibitem[Adadi and Berrada(2018)]{adadi2018xai}
Amina Adadi and Mohammed Berrada.
2018.
Peeking Inside the Black-Box: A Survey on Explainable Artificial Intelligence (XAI).
\emph{IEEE Access}, 6:52138--52160.

\bibitem[Agarwal et~al.(2022)Agarwal, Ley, Krishna, Saxena, Pawelczyk, Johnson, Puri, Zitnik, and Lakkaraju]{agarwal2022openxai}
Chirag Agarwal, Dan Ley, Satyapriya Krishna, Eshika Saxena, Martin Pawelczyk,
Nari Johnson, Isha Puri, Marinka Zitnik, and Himabindu Lakkaraju.
2022.
OpenXAI: Towards a Transparent Evaluation of Post hoc Model Explanations.
\emph{arXiv preprint arXiv:2206.11104}.

\bibitem[Agarwal et~al.(2023)Agarwal, Queen, Lakkaraju, and Zitnik]{agarwal2023gnn_eval}
Chirag Agarwal, Owen Queen, Himabindu Lakkaraju, and Marinka Zitnik.
2023.
Evaluating explainability for graph neural networks.
\emph{Scientific Data}, 10:144.

\bibitem[Ali et~al.(2023)Ali, Abuhmed, El-Sappagh, Muhammad, Alonso-Moral, Confalonieri, Guidotti, Del Ser, Diaz-Rodriguez, and Herrera]{ali2023trustworthy}
Safwan Ali, Tamer Abuhmed, Shaker El-Sappagh, Khan Muhammad, Jose M. Alonso-Moral,
Roberto Confalonieri, Riccardo Guidotti, Javier Del Ser, Natalia Diaz-Rodriguez,
and Francisco Herrera.
2023.
Explainable Artificial Intelligence (XAI): What we know and what is left to attain Trustworthy Artificial Intelligence.
\emph{Information Fusion}, 99:101805.

\bibitem[Arrieta et~al.(2020)]{arrieta2020xai}
Alejandro Barredo Arrieta, Natalia Diaz-Rodriguez, Javier Del Ser,
Adrien Bennetot, Siham Tabik, Alberto Barbado, Salvador Garcia,
Sergio Gil-Lopez, Daniel Molina, Richard Benjamins, Raja Chatila,
and Francisco Herrera.
2020.
Explainable Artificial Intelligence (XAI): Concepts, taxonomies, opportunities and challenges toward responsible AI.
\emph{Information Fusion}, 58:82--115.

\bibitem[Canha et~al.(2025)Canha, Kubler, Framling, and Fagherazzi]{canha2025functionally}
Dulce Canha, Sylvain Kubler, Kary Framling, and Guy Fagherazzi.
2025.
A Functionally-Grounded Benchmark Framework for XAI Methods: Insights and Foundations from a Systematic Literature Review.
\emph{ACM Computing Surveys}, 57(12):1--40.

\bibitem[Gu et~al.(2024)Gu, Jiang, Shi, Tan, Zhai, Xu, Li, Shen, Ma, Liu, Wang, Zhang, Wang, Gao, Ni, and Guo]{gu2024llmjudge}
Jiachen Gu, Xiaofeng Jiang, Zhen Shi, Han Tan, Xiaoyi Zhai, Chang Xu,
Wei Li, Yutong Shen, Shiyang Ma, Hongwei Liu, Sheng Wang, Kai Zhang,
Yifei Wang, Wei Gao, Lionel Ni, and Jun Guo.
2024.
A Survey on LLM-as-a-Judge.
\emph{arXiv preprint arXiv:2411.15594}.

\bibitem[Hedstrom et~al.(2023)Hedstrom, Weber, Bareeva, Krakowczyk, Motzkus, Samek, Lapuschkin, and Hohne]{hedstrom2023quantus}
Anna Hedstrom, Leander Weber, Dilyara Bareeva, Daniel Krakowczyk,
Franz Motzkus, Wojciech Samek, Sebastian Lapuschkin, and Marina M.-C. Hohne.
2023.
Quantus: An Explainable AI Toolkit for Responsible Evaluation of Neural Network Explanations and Beyond.
\emph{Journal of Machine Learning Research}, 24:1--11.

\bibitem[Kadir et~al.(2023)Kadir, Mosavi, and Sonntag]{kadir2023metrics}
Md Abdul Kadir, Amir Mosavi, and Daniel Sonntag.
2023.
Evaluation Metrics for XAI: A Review, Taxonomy, and Practical Applications.
In \emph{Proceedings of the 2023 IEEE 27th International Conference on Intelligent Engineering Systems}, 111--124.

\bibitem[Longo et~al.(2024)Longo, Brcic, Cabitza, Choi, Confalonieri, Del Ser, Guidotti, Hayashi, Herrera, Holzinger, Jiang, Khosravi, Lecue, Malgieri, Paez, Samek, Schneider, Speith, and Stumpf]{longo2024manifesto}
Luca Longo, Mario Brcic, Federico Cabitza, Jaeho Choi, Roberto Confalonieri,
Javier Del Ser, Riccardo Guidotti, Yuki Hayashi, Francisco Herrera,
Andreas Holzinger, Renzhe Jiang, Hassan Khosravi, Freddy Lecue,
Gianclaudio Malgieri, Angel Paez, Wojciech Samek, Johannes Schneider,
Timon Speith, and Simone Stumpf.
2024.
Explainable artificial intelligence (XAI) 2.0: A manifesto of open challenges and interdisciplinary research directions.
\emph{Information Fusion}, 106:102301.

\bibitem[Lundberg and Lee(2017)]{lundberg2017unified}
Scott M. Lundberg and Su-In Lee.
2017.
A Unified Approach to Interpreting Model Predictions.
In \emph{Advances in Neural Information Processing Systems}, volume 30,
4765--4774.

\bibitem[Mothilal et~al.(2020)Mothilal, Sharma, and Tan]{mothilal2020dice}
Ramaravind K. Mothilal, Amit Sharma, and Chenhao Tan.
2020.
Explaining machine learning classifiers through diverse counterfactual explanations.
In \emph{Proceedings of the 2020 Conference on Fairness, Accountability, and Transparency}, 607--617.

\bibitem[Pawlicki et~al.(2024)Pawlicki, Pawlicka, Uccello, Szelest, D'Antonio, Kozik, and Choras]{pawlicki2024multiple}
Marek Pawlicki, Aleksandra Pawlicka, Federica Uccello, Sebastian Szelest,
Salvatore D'Antonio, Rafal Kozik, and Michal Choras.
2024.
Evaluating the necessity of the multiple metrics for assessing explainable AI: A critical examination.
\emph{Neurocomputing}, 602:128282.

\bibitem[Proszewska et~al.(2025)Proszewska, Danel, and Rymarczyk]{proszewska2025bxaic}
Magdalena Proszewska, Tomasz Danel, and Dawid Rymarczyk.
2025.
B-XAIC Dataset: Benchmarking Explainable AI for Graph Neural Networks Using Chemical Data.
\emph{arXiv preprint arXiv:2505.22252}.

\bibitem[Retzlaff et~al.(2024)Retzlaff, Angerschmid, Saranti, Schneeberger, Rottger, Muller, and Holzinger]{retzlaff2024posthoc}
Christoph O. Retzlaff, Anna Angerschmid, Alexandra Saranti, Dominik Schneeberger,
Richard Rottger, Heimo Muller, and Andreas Holzinger.
2024.
Post-hoc vs ante-hoc explanations: xAI design guidelines for data scientists.
\emph{Cognitive Systems Research}, 86:101243.

\bibitem[Ribeiro et~al.(2016)Ribeiro, Singh, and Guestrin]{ribeiro2016why}
Marco Tulio Ribeiro, Sameer Singh, and Carlos Guestrin.
2016.
Why Should I Trust You?: Explaining the Predictions of Any Classifier.
In \emph{Proceedings of the 22nd ACM SIGKDD International Conference on Knowledge Discovery and Data Mining}, 1135--1144.

\bibitem[Rudin(2019)]{rudin2019stop}
Cynthia Rudin.
2019.
Stop explaining black box machine learning models for high stakes decisions and use interpretable models instead.
\emph{Nature Machine Intelligence}, 1(5):206--215.

\bibitem[Wu et~al.(2025)Wu, Kaushik, Li, Bauer, and Onoue]{wu2025userperceptions}
Xinru Wu, Rohan Kaushik, Wenxin Li, Lujo Bauer, and Kenta Onoue.
2025.
User perceptions vs. proxy LLM judges: Privacy and helpfulness in LLM responses to privacy-sensitive scenarios.
\emph{arXiv preprint arXiv:2510.20721}.

\bibitem[Zheng et~al.(2023)Zheng, Chiang, Sheng, Zhuang, Wu, Zhuang, Lin, Li, Li, Xing, Zhang, Gonzalez, and Stoica]{zheng2023judging}
Lianmin Zheng, Wei-Lin Chiang, Ying Sheng, Siyuan Zhuang, Zihao Wu,
Yuxin Zhuang, Zhiruo Lin, Zhuohan Li, Dacheng Li, Eric P. Xing,
Hao Zhang, Joseph E. Gonzalez, and Ion Stoica.
2023.
Judging LLM-as-a-Judge with MT-Bench and Chatbot Arena.
In \emph{Advances in Neural Information Processing Systems 36: Datasets and Benchmarks Track}.

\bibitem[Zheng et~al.(2025)Zheng, Shirani, Chen, Lin, Cheng, Guo, and Luo]{zheng2025ffidelity}
Xu Zheng, Farhad Shirani, Zhuomin Chen, Chaohao Lin, Wei Cheng,
Wenbo Guo, and Dongsheng Luo.
2025.
F-Fidelity: A Robust Framework for Faithfulness Evaluation of Explainable AI.
In \emph{International Conference on Learning Representations}.

\bibitem[Zhou et~al.(2025)Zhou, Xie, Li, Zhan, Song, Yang, Espinoza, Welton, Mai, Jin, Xu, Chung, Xing, Tsai, Schaffer, Shi, Liu, Liu, and Zhang]{zhou2025medthink}
Shuyi Zhou, Weijia Xie, Jie Li, Zijie Zhan, Mengyi Song, Haoyu Yang,
Carlos Espinoza, Lauren Welton, Xiang Mai, Yichen Jin, Zhenzhen Xu,
Yu-Han Chung, Yuhui Xing, Ming-Han Tsai, Eric Schaffer, Ya Shi,
Ning Liu, Zhiyong Liu, and Rui Zhang.
2025.
Automating expert-level medical reasoning evaluation of large language models.
\emph{npj Digital Medicine}, 9(1):34.

\end{thebibliography}

\end{document}
